\documentclass[10pt]{article}
\usepackage{amsmath, amssymb}
\usepackage{geometry}
\geometry{a4paper, margin=0.7in}
\usepackage{mathtools}
\usepackage{lscape}
\usepackage{graphicx}
\usepackage{enumitem}
\usepackage{caption}
\usepackage{multicol}
\usepackage{booktabs}
\usepackage{enumitem}

\usepackage{kotex}
\usepackage{fontspec}
\setmainfont{IBMPlexSans}[
    Path = ./IBMPlexSans/,
    Extension = .ttf,
    UprightFont = *-Regular,
    BoldFont = *-Bold,
    ItalicFont = *-Italic,
    BoldItalicFont = *-BoldItalic
]

\usepackage{amsthm}
\usepackage[most]{tcolorbox}

\makeatletter
\renewenvironment{proof}[1][\proofname]{%
  \par\vspace{1ex}\pushQED{\qed}\normalfont
  \topsep6\p@\@plus6\p@\relax
  \trivlist
  \item[\hskip\labelsep\bfseries #1\@addpunct{.}]%
}{%
  \popQED\endtrivlist\@endpefalse
  \addvspace{1ex}%
}
\makeatother

% ---- Custom style for unnumbered, clean theorem box ----
\tcbset{
  myboxstyle/.style={
    colback=gray!10,       % 회색 배경
    colframe=gray!60,      % 테두리 색
    boxrule=0.4pt,         % 테두리 두께
    arc=1.5mm,               % 모서리 둥글게
    left=6pt, right=6pt, top=6pt, bottom=6pt,
  }
}

\newtcolorbox{stabilitybox}[1]{myboxstyle, title=\textbf{#1}, fonttitle=\bfseries\color{black}}

\newtheorem{lemma}{Lemma}
\newtheorem{proposition}{Proposition}
\newtheorem{corollary}{Corollary}
\newtheorem{definition}{Definition}
\newtheorem{example}{Example}
\newtheorem{remark}{Remark}

\newcounter{remarkctr}
\renewenvironment{remark}[1][]{%
  \refstepcounter{remarkctr}% 카운터 증가
  \par\vspace{1ex}\noindent
  \textbf{Remark~\theremarkctr.}%
  \ifx#1\empty\else\ \textbf{#1}\fi\par
}{\vspace{1ex}}

\begin{document}

\begin{center}
    {\LARGE\bfseries IDO6001 Homework 03 \par}
    \vspace{1em}
    {\large 장서현 (학번: 2025324096) \par}
\end{center}
\vspace{0.5em}

% -------------------------------------------------------------------------------------
\section{Problem 1}
다음의 $(1||\sum_{i=1}^n T_i)$ 문제를 (1) 선형정수계획법으로 모델링하고, (2) 선형계획이완(LP relaxation) 이용한 하한경계를 사용하여 분지한계법으로 깊이 탐색과 너비 탐색 모두 사용하여 푸시오. 
\begin{table}[h!]
    \centering
    \begin{tabular}{c|ccccc}
    \hline
    작업 & 1 & 2 & 3 & 4 & 5 \\ \hline
    가공시간 & 4 & 3 & 7 & 2 & 2 \\
    납기 & 5 & 6 & 8 & 8 & 17 \\ \hline
    \end{tabular}
    \end{table}

\subsection{선형정수계획법 모델}
\begin{stabilitybox}{Notation}
    \setlist[itemize]{leftmargin=2em, itemsep=1.5pt}
    \begin{minipage}[t]{0.51\linewidth}
        \textbf{Set and Indices}
        \begin{itemize}
            \item $N = \{1,2,3,4,5\}$: 작업의 집합
            \item $i \in N$: 작업의 인덱스
            \item $T = \{1, \dots, |T|\}$: 시점의 집합
            \item $t, s \in T$: 시점의 인덱스
            \begin{itemize}
                \item 각 시점 t는 시간구간 $[t-1, t)$의 시작 시점을 의미
            \end{itemize}
        \end{itemize}
    \end{minipage}
    \hfill
    \begin{minipage}[t]{0.48\linewidth}
        \textbf{Decision Variables}
        \begin{itemize}
            \item $x_{it}$: 작업 $i$가 시점 $t$에 시작하면 1, 아니면 0
        \end{itemize}
        \textbf{Parameters}
        \begin{itemize}
            \item $p_i$: 작업 $i$의 가공 시간
            \item $d_i$: 작업 $i$의 납기
            \item $|T| = \sum p_i$: 최대 허용 시간
            \item $\xi_{it}$: 작업 $i$를 시점 $t$에 시작할 경우 발생하는 지연시간
            \begin{itemize}
                \item $\xi_{it} = \max(0, t-1 + p_i - d_i) \quad \forall i \in N, t \in T$
            \end{itemize}
        \end{itemize}
    \end{minipage}
\end{stabilitybox}

\subsubsection*{Model Formulation using Time-Indexed Variables}
\begin{subequations}
    \begin{align}
        \text{min} \quad & \sum_{i \in N} \sum_{t=1}^{|T|-p_i+1} \xi_{it}\, x_{it} \label{1a}\\
        \text{s.t.} \quad & \sum_{t=1}^{|T|-p_i+1} x_{it} = 1  & \forall i \in N \label{1b}\\
        & \sum_{i \in N} \sum_{s=\max(1, t-p_i+1)}^{t} x_{is} \le 1 & \forall t \in T \label{1c}\\
        & x_{it} \in \{0, 1\} & \forall i \in N, t \in T \label{1d}
    \end{align}
\end{subequations}
\vspace{0.05em}
\begin{itemize} [leftmargin=1em, itemsep=0pt, topsep=1pt]
    \item \eqref{1a} 총 지연 시간의 합을 최소화하는 목적함수
    \item \eqref{1b} 모든 작업 $i$는 스케줄링 기간($T$) 내에 한 번만 시작되어야 함 
    \item \eqref{1c} 모든 시점 $t$에 기계는 최대 1개의 작업만 처리할 수 있음
    \item \eqref{1d} $x_{it}$가 시작 시점을 결정하는 이진 변수임을 정의
    \end{itemize}
\subsubsection*{LP relaxed Model Formulation}
\eqref{1a} $\sim$ \eqref{1c}는 동일하게 유지되며, \eqref{1d}가 아래의 제약식 \eqref{1e}으로 변환
\begin{equation*}
    0 \le x_{it} \le 1 \qquad \forall i \in N, t \in T \tag{1e} \label{1e}
\end{equation*}

\newpage

\begin{stabilitybox}{Approach}
    \begin{itemize}[leftmargin=1em]
        \item 분지 조건: 현재 노드의 LP Relaxation을 푼 해에 $0 < x_{it} < 1$인 변수가 하나라도 존재할 때,
        \begin{enumerate}
            \item 실수 값을 갖는 모든 $x_{it}$ 중에서, 시점 $t$가 가장 빠른 $x_{it}$
            \item 규칙 1에서 동일한 시점 $t$에 대해 변수가 여러 개 존재할 시, 작업 $i$의 납기가 가장 빠른 $x_{it}$
        \end{enumerate}
        \item 하한경계: 각 노드에서 추가된 제약을 포함한 LP Relaxation의 목적함수 최적값
        \item 가지치기 조건: 
        \begin{enumerate}
            \item Bound: 현재 노드의 LB $\ge$ IC
            \item Integral: 현재 노드의 LB 해가 모든 $x_{it}$에 대해 정수해
            \item Infeasible: 현재 노드의 LP Relaxation 문제가 infeasible
        \end{enumerate}
        \item 초기 IC = $|T|$ = 18 
    \end{itemize}
\end{stabilitybox}

\subsection{깊이탐색을 활용한 분지한계법}

\subsection{너비탐색을 활용한 분지한계법} 

% -------------------------------------------------------------------------------------
\newpage
\section{Problem 2}
다음의 $(1|s_{ij}|C_{max})$ 문제를 다음의 작업준비시간을 고려하여 (1) 선형정수계획법으로 모델링하고, (2) 선형계획이완(LP relaxation) 이용한 하한경계를 사용하여 분지한계법으로 깊이 탐색과 너비 탐색 모두 사용하여 푸시오. 
\begin{table}[h!]
\centering
\begin{tabular}{c|cccc}
\hline
 & 1 & 2 & 3 & 4 \\ \hline
1 & -- & 30 & 50 & 90 \\
2 & 40 & -- & 20 & 80 \\
3 & 30 & 30 & -- & 60 \\
4 & 20 & 15 & 10 & -- \\ \hline
\end{tabular}
\end{table}

\subsection{선형정수계획법 모델}
\begin{stabilitybox}{Notation}
    \setlist[itemize]{leftmargin=2em, itemsep=1.5pt}
    \begin{minipage}[t]{0.58\linewidth}
        \textbf{Set and Indices}
        \begin{itemize}
            \item $N = \{1, 2, 3, 4\}$: 실제 작업 집합
            \item $N_0 = \{0, 1, 2, 3, 4\}$: 가상 작업을 포함한 전체 작업 집합
            \item $i, j \in N_0$: 노드를 나타내는 인덱스
        \end{itemize}
        \textbf{Decision Variables}
        \begin{itemize}
            \item $x_{ij}$: 작업 $i$가 작업 $j$ 바로 직전에 수행되면 1, 아니면 0 ($i \ne j$)
            \item $u_i$: 작업 $i$의 순서를 나타내는 보조 변수 ($i \in N$)
        \end{itemize}
    \end{minipage}
    \hfill
    \begin{minipage}[t]{0.4\linewidth}
        \textbf{Parameters}
        \begin{itemize}
            \item $s_{ij}$: 작업 $i$에서 작업 $j$로의 작업 준비 시간
            \begin{itemize}
                \item $s_{0j} = 0 \quad (\forall j \in N)$
                \item $s_{i0} = 0 \quad (\forall i \in N)$
            \end{itemize}
            \item $n = |N| = 4$: 실제 작업의 수
        \end{itemize}
    \end{minipage}
\end{stabilitybox}

\subsubsection*{Model Formulation using TSP Framework}
\begin{subequations}
    \begin{align}
        \text{min} \quad & \sum_{i \in N_0} \sum_{j \in N_0, j \ne i} s_{ij} x_{ij} \label{2a}\\
        \text{s.t.} \quad
        & \sum_{j \in N_0, j \ne i} x_{ij} = 1 & \forall i \in N_0 \label{2b}\\
        & \sum_{i \in N_0, i \ne j} x_{ij} = 1 & \forall j \in N_0 \label{2c}\\
        & u_i - u_j + n x_{ij} \le n - 1 & \forall i, j \in N, i \ne j \label{2d}\\
        & 1 \le u_i \le n & \forall i \in N \label{2e} \\
        & x_{ij} \in \{0, 1\} & \forall i, j \in N_0, i \ne j \label{2f}
    \end{align}
\end{subequations}
\vspace{0.05em}
\begin{itemize} [leftmargin=1em, itemsep=0pt, topsep=1pt]
    \item \eqref{2a} $C_{max}$는 모든 작업의 총 가공 시간과 모든 총 작업 준비 시간($\sum s_{ij}$)의 합이며, 총 가공 시간은 상수이기 떄문에, 목적함수는 총 작업 준비 시간을 최소화하는 것으로 설계
    \item \eqref{2b} 모든 작업 $i \in N_0$는, 스케줄 상에서 자신 바로 뒤에 위치할 단 하나의 작업 $j \in N_0$를 가져야 함
    \item \eqref{2c} 모든 작업 $j \in N_0$는, 스케줄 상에서 자신 바로 앞에 위치할 단 하나의 작업 $i \in N_0$를 가져야 함
    \item \eqref{2d} 서브투어 제거 제약: $x_{ij}=1$일 때 $u_i + 1 \le u_j$를 강제하여 15, 경로를 따라 $u$ 값이 증가하게 함으로써 $N$ 내의 닫힌 서브투어를 방지
    \item \eqref{2e} 스케줄에서 작업 $i$의 순서 $u_i$의 범위 정의
    \item \eqref{2f} $x_{ij}$가 이진 변수임을 정의
\end{itemize}

\subsubsection*{LP relaxed Model Formulation}
\eqref{2a} $\sim$ \eqref{2e}는 동일하게 유지되며, \eqref{2f}가 아래의 제약식 \eqref{2g}으로 변환
\begin{equation*}
    0 \le x_{ij} \le 1 \qquad \forall i, j \in N_0, i \ne j \tag{2g} \label{2g}
\end{equation*}

\newpage
\begin{stabilitybox}{Approach}
    \begin{itemize}[leftmargin=1em]
        \item 분지 조건: 현재 노드의 LP Relaxation을 푼 해에 $0 < x_{it} < 1$인 변수가 하나라도 존재할 때,
        \begin{enumerate}
            \item 실수 값을 갖는 모든 $x_{ij}$ 변수 중에서, 그 값이 0.5에 가장 가까운 $x_{ij}$
            \item 규칙 1에서 동률(tie)이 발생할 경우, 준비 시간($s_{ij}$)이 가장 작은 $x_{ij}$
        \end{enumerate}
        \item 하한경계: 각 노드에서 추가된 제약을 포함한 LP Relaxation의 목적함수 최적값
        \item 가지치기 조건: 
        \begin{enumerate}
            \item Bound: 현재 노드의 LB $\ge$ IC
            \item Integral: 현재 노드의 LB 해가 모든 $x_{it}$에 대해 정수해
            \item Infeasible: 현재 노드의 LP Relaxation 문제가 infeasible
        \end{enumerate}
        \item 초기 IC = 70: CUT 기반 설정 (4 $\to$ 3 $\to$ 1 $\to$ 2)
    \end{itemize}
\end{stabilitybox}

\subsection{깊이탐색을 활용한 분지한계법}
분지조건, 하한경계 명시


\subsection{너비탐색을 활용한 분지한계법} 
분지조건, 하한경계 명시


% -------------------------------------------------------------------------------------
\newpage
\section*{Appendix - Code Implementation}

    
\end{document}

% \begin{stabilitybox}{Notation}
%     \setlist[itemize]{leftmargin=2em, itemsep=1.5pt}
%     \begin{minipage}[t]{0.4\linewidth}
%         \textbf{Set and Indices}
%         \begin{itemize}
%             \item $N = \{1, 2, 3, 4, 5\}$: 작업의 집합
%             \item $j, k \in N$: 작업을 나타내는 인덱스
%         \end{itemize}
%         \textbf{Parameters}
%         \begin{itemize}
%             \item $p_j$: 작업 $j$의 가공 시간
%             \item $d_j$: 작업 $j$의 납기 
%             \item $M = 2\sum_{j\in N} p_j$: 양의 상수
%         \end{itemize}
%     \end{minipage}
%     \hfill
%     \begin{minipage}[t]{0.58\linewidth}
%         \textbf{Decision Variables}
%         \begin{itemize}
%             \item $C_j$: 작업 $j$의 완료 시간
%             \item $T_j$: 작업 $j$의 지연 시간
%             \item $y_{jk}$: 작업 $j$가 작업 $k$보다 먼저 처리되면 1, 아니면 0 ($j < k$)
%         \end{itemize}
%     \end{minipage}
% \end{stabilitybox}

% \subsubsection*{Model Formulation using Time-Indexed Variables}
% \begin{subequations}
%     \begin{align}
%         \text{min} \quad & \sum_{j \in N} T_j \label{1a}\\
%         \text{s.t.} \quad
%         & C_j \ge p_j & \forall j \in N \label{1b}\\
%         & C_j + p_k \le C_k + M(1 - y_{jk}) & \forall j, k \in N, j < k \label{1c}\\
%         & C_k + p_j \le C_j + M y_{jk} & \forall j, k \in N, j < k \label{1d}\\
%         & T_j \ge C_j - d_j & \forall j \in N \label{1e}\\
%         & C_j, T_j \ge 0 & \forall j \in N \label{1f}\\
%         & y_{jk} \in \{0, 1\} & \forall j, k \in N, j < k \label{1g}
%     \end{align}
% \end{subequations}
% \vspace{0.05em}
% \begin{itemize} [leftmargin=1em, itemsep=0pt, topsep=1pt]
%     \item \eqref{1a} 총 지연 시간의 합을 최소화하는 목적 함수
%     \item \eqref{1b} 작업 $j$의 완료 시간은 최소한 자신의 가공 시간($p_j$)보다 커야 함을 보장
%     \item \eqref{1c} \eqref{1d} 두 작업 $j, k$의 순서를 강제하는 disjunctive constraints
%     \begin{itemize}[itemsep=0pt, topsep=1pt]
%         \item $y_{jk}=1$ ($j \to k$)이면, \eqref{1c}가 $C_j + p_k \le C_k$가 되어 $k$는 $j$가 끝난 후에만 완료될 수 있음을 강제
%         \item $y_{jk}=0$ ($k \to j$)이면, \eqref{1d}가 $C_k + p_j \le C_j$가 되어 $j$는 $k$가 끝난 후에만 완료될 수 있음을 강제
%     \end{itemize}
%     \item \eqref{1e} 지연 시간 $T_j$를 $T_j = \max(0, C_j - d_j)$로 정의하는 제약 중 일부
%     \item \eqref{1f} 완료 시간과 지연 시간은 음수가 될 수 없음
%     \item \eqref{1g} $y_{jk}$가 순서를 결정하는 이진 변수임을 정의
% \end{itemize}

% \subsubsection*{LP relaxed Model Formulation}
% \eqref{1a} $\sim$ \eqref{1f}는 동일하게 유지되며, \eqref{1g}가 아래의 제약식 \eqref{1h}으로 변환
% \begin{equation*}
%     0 \le y_{jk} \le 1 \qquad \forall j, k \in N, j < k \tag{1h} \label{1h}
% \end{equation*}